\section{数据扩展：从稀缺轨迹到可持续训练闭环}

\subsection{为什么数据是瓶颈}
Web Agent 的高质量交互数据比纯文本语料更昂贵。每条轨迹都包含环境状态、动作序列、工具结果和终局判定，采集与清洗成本都高。

\begin{importantbox}{数据难点}
\begin{itemize}
    \item 真实交互慢，单位样本成本高；
    \item 轨迹长，标注和验证复杂；
    \item 环境随时间变化，旧数据会过时。
\end{itemize}
\end{importantbox}

\subsection{在线收集与过滤}
讲者建议在真实或高保真模拟环境中在线采样，再通过 verifier/Judge 做分层过滤：先去除明显失败，再保留高置信度成功轨迹，最后按任务类型均衡采样。

\begin{knowledgebox}{可执行的数据流水线}
\begin{enumerate}
    \item 采样：多策略并行生成轨迹；
    \item 过滤：终局判定 + 过程判定；
    \item 重加权：困难样本提升权重；
    \item 训练：SFT 或 RL 更新；
    \item 回放：在线 A/B 检验是否真提升。
\end{enumerate}
\end{knowledgebox}

\subsection{训练阶段的组合策略}
实践中常见组合是 ``SFT 打底 + 搜索增强 + RL 微调''。SFT 提供稳定起点，搜索提供高质量探索轨迹，RL 在交互反馈下做策略细化。

\begin{table}[H]
\centering
\small
\begin{tabular}{p{2.4cm}p{3.9cm}p{4.7cm}}
\toprule
\textbf{阶段} & \textbf{目标} & \textbf{注意事项} \\
\midrule
SFT & 学会基本动作语法与任务格式 & 防止只记模板，忽视长链路泛化 \\
Search-augmented & 提升探索质量与成功轨迹占比 & 控制推理预算，避免过度采样 \\
RL / policy update & 对齐终局成功率和鲁棒性 & 防止 reward hacking 与分布偏移 \\
\bottomrule
\end{tabular}
\caption{Web Agent 常见训练组合路径}
\end{table}

\begin{warningbox}{数据扩展中的隐性风险}
如果只追求成功率，系统可能偏向 ``容易任务''，导致困难场景性能退化。数据扩展必须和任务难度分布管理同时进行。
\end{warningbox}

\subsection{本章小结}
Agent 能力扩展的核心不只是模型升级，而是可持续数据闭环。在线采样、可靠过滤和难度均衡决定了系统是否能稳定进步。

